\section{Discussion and conclusions}

The central positive result is modest but exact: a quantum process may create distinguishable internal records, compute a property from them, convert that property to phase, and uncompute every controlled record while retaining an interferometric effect. Internal processing and conditional action do not introduce a new fundamental source of decoherence. What matters is whether history-dependent information remains outside the inverse and whether accumulated control errors preserve the required phase.

The central negative result is equally important: the same endpoint channel can be implemented without the advertised internal history. Therefore, a final interference fringe cannot certify an agent, a memory depth, or a unique temporal narrative. The direct-phase control is supposed to agree. Historical authenticity becomes a scientifically distinct question only after an access model introduces intervention times, late inputs, locality, black-box dynamics, or resource bounds.

Zero record gives a useful accounting rule. It is not enough to set a named memory qubit to zero; every residual system must be checked for distinctions that the allowed predicate does not reveal. Conditional mutual information captures this within-predicate requirement, provided the predicate is non-injective. The four-history benchmark exists specifically to prevent a formally correct but vacuous privacy score.

The analytic late-challenge protocol moves one step beyond a static ancilla. A classical response at a later time depends on information acquired before the challenge, and a bounded classical-memory comparison gives a quantitative success limit. The coherent phase variant uses the same response function, but its endpoint task is not the classical random-access score, and the current simulator does not model challenge independence. Both statements remain architecture-relative. If a shortcut device can access the history label when the response is produced, it can insert the phase directly. A publishable causal advance requires a witness or interactive task that makes the forbidden access operational rather than rhetorical.

Three research directions are defensible. First, process-tensor and temporal-correlation tools may yield a targeted memory-dimension witness compatible with a later inverse. Second, quantum-algorithmic-measurement methods may establish a coherent-versus-incoherent query separation for a black-box history property. Third, simulator-based quantum zero-knowledge ideas may inspire a composable definition in which a verifier learns a predicate while any allowed residual view is simulable from that predicate alone. None of these results is claimed here.

The programme should be discontinued as a distinct research framing if every candidate task compiles to a known phase oracle under the same access assumptions, if memory cannot be causally witnessed, or if no new resource bound or experimentally distinguishable prediction results. A rigorous negative result would still be valuable: it would identify exactly why reversible internal narrative adds no operational resource beyond coherent computation.

In its present form, the project is a completed baseline rather than a completed research field. It supplies definitions, no-go statements, reference circuits, controls, reproducible numerics, and a falsifiable next target. Its scientific claim is not that quantum mechanics exposes alternative universes. It is that the boundary between a record-rich internal process and its recoverable coherence can be formulated as a precise causal-information task, and that any stronger claim must survive endpoint equivalence, transcript accounting, and explicit comparison models.
